\documentclass{szb-solution}

\area{Kalkulus}
\title{Differenciálszámítás III}
\subject{Matematika G1}
\subjectCode{BMETE94BG01}
\date{Utoljára frissítve: \today}
\docno{9}

\begin{document}
\maketitle

\subsection{Teljes függvényvizsgálat}

Legyen
$$
  f(x)=\frac{2x^2}{x^2-9}
  \text.
$$
Az értelmezési tartomány
$$
  \Domain_f=\mathbb R\setminus\{-3,3\}
  \text.
$$
A függvény páros, mert $f(-x)=f(x)$, ezért a grafikon az $y$ tengelyre
szimmetrikus. Egyetlen zérushelye $x=0$, és $f(0)=0$.

A nevezetes határértékek:
\begin{align*}
  \lim_{x\to\pm\infty}f(x)&=2,\\
  \lim_{x\to-3^-}f(x)&=+\infty,
  &\lim_{x\to-3^+}f(x)&=-\infty,\\
  \lim_{x\to3^-}f(x)&=-\infty,
  &\lim_{x\to3^+}f(x)&=+\infty
  \text.
\end{align*}
Így $x=-3$ és $x=3$ függőleges, $y=2$ pedig vízszintes
aszimptota.

Az első derivált
$$
  f'(x)=\frac{4x(x^2-9)-4x^3}{(x^2-9)^2}
       =\frac{-36x}{(x^2-9)^2}
  \text.
$$
A nevező pozitív, ezért a derivált előjelét $-x$ előjele dönti el. A
függvény szigorúan növekvő a $(-\infty,-3)$ és $(-3,0)$ intervallumokon,
szigorúan csökkenő a $(0,3)$ és $(3,\infty)$ intervallumokon. Az
$x=0$ pontban lokális maximuma van:
$$
  \boxed{f(0)=0}
  \text.
$$

A második derivált
$$
  f''(x)=\frac{108(x^2+3)}{(x^2-9)^3}
  \text.
$$
Mivel a számláló mindig pozitív, $f$ konkáv a $(-3,3)$ intervallumon,
és konvex a $(-\infty,-3)$, illetve $(3,\infty)$ intervallumon. A
konvexitás csak a nem értelmezett $\pm3$ pontokban változik, ezért nincs
inflexiós pont.

Az értékkészlethez oldjuk meg az $y=f(x)$ egyenletet $x^2$-re:
$$
  y=\frac{2x^2}{x^2-9}
  \quad\Longleftrightarrow\quad
  x^2=\frac{9y}{y-2}
  \text.
$$
Valós $x$ pontosan akkor létezik, ha a jobb oldal nemnegatív. Ezért
$$
  \boxed{\Range_f=(-\infty,0]\cup(2,\infty)}
  \text.
$$

\subsection{Legkisebb felszínű, felül nyitott henger}

Egy liter $1\,\mathrm{dm}^3$. Ha a henger sugara $r>0$, magassága $h>0$,
akkor
$$
  \pi r^2h=1,
  \qquad
  h=\frac{1}{\pi r^2}
  \text.
$$
A felül nyitott henger felszíne a palást és az alsó körlap összege:
$$
  A(r)=2\pi rh+\pi r^2=\frac{2}{r}+\pi r^2
  \text.
$$
Deriválva
$$
  A'(r)=-\frac{2}{r^2}+2\pi r
  \text,
$$
így $A'(r)=0$ pontosan akkor, ha $r^3=1/\pi$. Továbbá
$$
  A''(r)=\frac{4}{r^3}+2\pi>0
  \qquad(r>0),
$$
tehát ez valóban az egyetlen globális minimum. A hozzá tartozó magasság
megegyezik a sugárral:
$$
  \boxed{r=h=\pi^{-1/3}\,\mathrm{dm}}
  \text,
  \qquad
  \boxed{A_{\min}=3\pi^{1/3}\,\mathrm{dm}^2}
  \text.
$$

\subsection{Legnagyobb térfogatú, adott alkotójú kúp}

Jelölje az adott alkotó hosszát $h$, az alkotó és a tengely által bezárt
szöget pedig $\alpha\in(0,\pi/2)$. Ekkor a kúp alapkörének sugara és
magassága
$$
  r=h\sin\alpha,
  \qquad
  m=h\cos\alpha
  \text.
$$
A térfogat tehát
$$
  V(\alpha)=\frac{\pi h^3}{3}\sin^2\alpha\cos\alpha
  \text.
$$
Mivel a pozitív állandó szorzó a szélsőérték helyét nem befolyásolja,
elég a $\sin^2\alpha\cos\alpha$ kifejezést vizsgálni. Deriváltja
$$
  \sin\alpha\bigl(2\cos^2\alpha-\sin^2\alpha\bigr)
  \text.
$$
A belső kritikus pontban
$$
  2\cos^2\alpha=\sin^2\alpha,
  \qquad
  \sin^2\alpha=\frac23,
  \qquad
  \cos\alpha=\frac1{\sqrt3}
  \text.
$$
A térfogat az intervallum mindkét végén nullához tart, ezért ez a pont adja
a maximumot. Így
$$
  \boxed{r=h\sqrt{\frac23},\qquad m=\frac{h}{\sqrt3}}
$$
és
$$
  \boxed{V_{\max}=\frac{2\pi h^3}{9\sqrt3}}
  \text.
$$

\subsection{Körbe írt legnagyobb területű téglalap}

Legyen a téglalap két oldalhossza $a,b>0$. A téglalap átlója a kör
átmérője, ezért
$$
  a^2+b^2=4r^2
  \text.
$$
Az $a^2+b^2$ rögzített összeg mellett
$$
  2ab\leq a^2+b^2=4r^2
  \text,
$$
és egyenlőség pontosan $a=b$ esetén áll fenn. A maximális területű
téglalap tehát négyzet, melynek
$$
  \boxed{a=b=\sqrt2\,r}
  \qquad\text{és}\qquad
  \boxed{T_{\max}=2r^2}
  \text.
$$

\subsection{A folyosósaroknál befordítható gerenda}

Legyen $\alpha\in(0,\pi/2)$ a gerenda és a $b$ szélességű folyosó
hossziránya által bezárt szög. A sarkot éppen érintő gerenda két része
$$
  \ell_1=\frac{a}{\sin\alpha},
  \qquad
  \ell_2=\frac{b}{\cos\alpha}
  \text,
$$
így az akadályt jelentő teljes hossz
$$
  \ell(\alpha)=\frac{a}{\sin\alpha}+\frac{b}{\cos\alpha}
  \text.
$$
A befordítható legnagyobb gerendahossz e mennyiség minimuma. Deriválva
$$
  \ell'(\alpha)
  =-\frac{a\cos\alpha}{\sin^2\alpha}
   +\frac{b\sin\alpha}{\cos^2\alpha}
  \text.
$$
A kritikus pont feltétele
$$
  b\sin^3\alpha=a\cos^3\alpha,
  \qquad
  \tan\alpha=\sqrt[3]{\frac ab}
  \text.
$$
Mivel $\ell(\alpha)$ mindkét végponton $+\infty$-hez tart, a belső
kritikus pont globális minimum. Behelyettesítés után
$$
  \boxed{\ell_{\max}
  =\left(a^{2/3}+b^{2/3}\right)^{3/2}}
  \text.
$$

\subsection{A gazda legrövidebb útja}

Feltesszük, hogy a szárazföldön mindenütt azonos sebességgel halad, és a
pataknál elegendő egy pontot érintenie. Jelölje $x$ kilométerben a ház
patakra eső merőleges vetülete és az érintési pont közötti távolságot a
kocsma irányában. Ekkor a teljes út
$$
  L(x)=\sqrt{x^2+0{,}25^2}
       +\sqrt{(1-x)^2+0{,}75^2},
  \qquad 0\leq x\leq1
  \text.
$$
Deriváltja
$$
  L'(x)=\frac{x}{\sqrt{x^2+0{,}25^2}}
       -\frac{1-x}{\sqrt{(1-x)^2+0{,}75^2}}
  \text.
$$
Az $L'(x)=0$ egyenletben mindkét oldal nemnegatív, ezért négyzetre emelhetünk.
Egyszerűsítés után
$$
  0{,}75^2x^2=0{,}25^2(1-x)^2,
  \qquad
  0{,}75x=0{,}25(1-x),
$$
tehát
$$
  \boxed{x=0{,}25\,\mathrm{km}}
  \text.
$$
Ezt a visszatükrözés módszere is igazolja: a házat a patakra tükrözve a
legrövidebb törtvonal kiegyenesedik. A minimális teljes út hossza
$$
  L_{\min}
  =\sqrt{0{,}25^2+0{,}25^2}
   +\sqrt{0{,}75^2+0{,}75^2}
  =\boxed{\sqrt2\,\mathrm{km}}
  \text.
$$

\end{document}
